%% =====================================================================
%%  X-UAD : A Unified, Uncertainty-Aware and Self-Explaining Detector for
%%          Physical and Digital Face Attacks   (revision 2)
%%  IEEE conference format (IEEEtran).  Compile with pdflatex + bibtex:
%%     pdflatex xuad_v2 && bibtex xuad_v2 && pdflatex xuad_v2 && pdflatex xuad_v2
%%  Prose follows the revised draft; figures and supplementary tables are
%%  restored from xuad.tex; every number is \input from tables/ (generated by
%%  scripts/make_tables.py from results/*.json).
%% =====================================================================
\documentclass[conference]{IEEEtran}
\IEEEoverridecommandlockouts

\usepackage{cite}
\usepackage{amsmath,amssymb,amsfonts}
\usepackage{graphicx}
\usepackage{booktabs}
\usepackage{multirow}
\usepackage{xcolor}
\usepackage{textcomp}
\usepackage{url}
\usepackage[hidelinks]{hyperref}

\newcommand{\bs}[1]{\boldsymbol{#1}}
\newcommand{\R}{\mathbb{R}}
\newcommand{\ket}[1]{\lvert #1 \rangle}
\newcommand{\bra}[1]{\langle #1 \rvert}
\newcommand{\braket}[2]{\langle #1 \vert #2 \rangle}

\begin{document}

\title{X-UAD: A Unified, Uncertainty-Aware and Self-Explaining Detector for
Physical and Digital Face Attacks}

\author{}

\maketitle

\begin{abstract}
Face recognition systems can be defeated by presentation attacks staged in front
of the camera, such as a printed photograph, a replayed video or a silicone mask,
and by digital manipulation of the captured image through face swapping,
expression reenactment, synthesis or adversarial perturbation. A detector built for
one family is usually blind to the other, and the recent unified detectors that
handle both still share three practical weaknesses. They emit a single uncalibrated
score with no way of signalling that an input resembles nothing they were trained
on; they cannot point to where in the frame the evidence sits; and when they do
explain themselves, the explanation comes from a language model that never observed
the detector's internal state. We present X-UAD, one model aimed at all three.
Adaptive Quantum Attention adds a Born-rule scoring channel beside ordinary
dot-product attention. It is sign-invariant by construction, which matches the
interference character of moir\'e, halftone and checkerboard evidence, and it costs
$\mathcal{O}(dL)$ parameters rather than $\mathcal{O}(d^{2})$. The same learnable
rotations define an orthonormal evidence basis whose measurements serve
simultaneously as classification features and, completely unchanged, as the
explanation shown to a reader. A decoupled evidential head predicts class direction
and evidence magnitude separately, and pseudo-unseen evidence suppression trains
that magnitude on blends of two attack families, an outlier proxy built from the
attack taxonomy itself. On a difficulty-calibrated, identity-disjoint benchmark of
eight attack families with pixel-level ground truth, X-UAD reaches $0.42\%$ ACER on
seen attacks. That is statistically level with the strongest baseline, CDCN++
($0.46\%$), and well ahead of ResNet-18 ($1.88\%$), although X-UAD trails both on
eight-way family accuracy ($94.1\%$ against $95.2\%$ and $96.9\%$). It localises the
evidence best, at $80.1\%$ region-level IoU over the full test set against $73.2\%$ for CDCN++. On attack
families withheld from training it has the highest average detection rate
($77.5\%$), and the synthetic-blend objective is worth $5.4$ AUROC points on the
uncertainty channel. The results have clear limits: every method collapses under a
change of capture domain, a white-box attacker can suppress X-UAD's uncertainty
without changing its decision, and all numbers come from a single seed on a
simulated benchmark. Code, the benchmark generator and configuration files mapping
the same pipeline onto public capture-based datasets are released.
\end{abstract}

\begin{IEEEkeywords}
face anti-spoofing, presentation attack detection, deepfake detection, unified
attack detection, vision transformer, quantum-inspired attention, evidential deep
learning, uncertainty estimation, open-set recognition, explainable AI
\end{IEEEkeywords}

% =====================================================================
\section{Introduction}
% =====================================================================
Consider everything that can go wrong in front of a face-unlock camera. Someone
could hold up a printed photo of you. They could replay a video of your face on a
phone screen. They could wear a silicone mask shaped like your face. These are
\emph{physical} attacks, and each leaves a physical fingerprint behind: the halftone
dot pattern of a printer, the pixel grid of a phone screen aliasing into a moir\'e
pattern, the waxy sheen and flattened texture of silicone.

Now consider what can go wrong \emph{after} the camera has captured a real,
unmanipulated photograph. Someone could swap in a different identity. They could
drive your expression with someone else's performance. They could generate a face
that never belonged to any real person. They could add a barely visible pattern of
noise that flips a classifier's decision without a human noticing anything at all.
These are \emph{digital} attacks, and their fingerprints look entirely different:
blending seams where a swapped face meets the original, resolution mismatches,
checkerboard patterns left behind by the upsampling layers of a generator, unusual
peaks in the frequency spectrum.

Because these two sets of fingerprints look nothing alike, the two research
communities that study them grew up separately. Face anti-spoofing work focused on
texture and depth cues \cite{liu2018auxiliary,yu2020cdcn,yu2021revisiting,wang2022ssan}.
Face forgery research focused on frequency artifacts and blending cues
\cite{rossler2019ff,li2020celebdf,frank2020frequency,wang2020cnnspot}. The split is
understandable, but it is not something you can deploy. An operator needs one
decision, one threshold and one audit trail, not two models that might disagree
with no principled way to combine them. A recent line of work
\cite{yu2024unified,yuan2024challenge,ma2025uniattackplus,chen2025uniattack,zhang2024unidefense}
has finally built unified corpora and joint detectors that handle both attack
families at once, which is real progress. It also exposed a problem that had been
hiding underneath. On the attack types these models were trained on they are
essentially saturated, with ACER around $0.5$ to $1.4\%$, but the moment you hold out
an attack family during training and show it only at test time, error rates jump to
between $22$ and $44\%$ with large variance between runs \cite{yu2024unified}.
Unified detection is close to solved in-distribution and still wide open
out-of-distribution.

Layered on top of that gap, we see three capabilities missing from almost all of
this literature, and an operator needs all three the moment a
score comes back ambiguous rather than obviously fake or obviously real.

The first is a real notion of \emph{not knowing}. A softmax output is fundamentally
a closed-set quantity, since it has to redistribute all of its probability mass
across the classes it was trained on. A completely novel kind of attack therefore
gets assigned high confidence toward whichever known class it happens to resemble
most. What is needed is a second-order belief, one that is allowed to sit
close to empty when the input does not look like anything the model has learned.

The second is \emph{evidence localisation}. Depth-map supervision
\cite{liu2018auxiliary,yu2021revisiting} works reasonably for physical attacks,
since a printed photo or a replayed video really is flatter than a real face. But
depth is a proxy for liveness rather than a map of forensic evidence, and it says
nothing useful about a digitally swapped face that came from a perfectly live,
perfectly three-dimensional original capture.

The third is an \emph{explanation that is tied to the decision}. Recent
work has started feeding face images into multimodal language models and asking
them to describe what looks fake \cite{huang2024ffaa,guo2025m2f2det}. The results
read fluently, but the language model never sees the detector's internal state. It
is reasoning over the raw image, independently of whatever the classifier
computed. Its narrative is therefore correlated with the decision at best, and it
is free to mention artifacts such as ``a visible screen bezel'' that are not
present in the frame or that the detector never used.

\subsection{What this paper does}
We built X-UAD, a single model that detects physical and digital face attacks
together, reports how much it does not know, points to where the evidence is, and
explains its own decision in language that traces back to numbers it
computed. Three ideas make this work.

\textbf{Adaptive Quantum Attention (AQA).} Ordinary attention scores a pair of
tokens with an unbounded linear functional, the dot product. A great deal of the
evidence for face attacks, though, is fundamentally an interference phenomenon: a
moir\'e beat, a checkerboard pattern, a spectral peak. These show up in the
\emph{squared magnitude} of a projection, not in the raw projection, and their sign
carries no information. A moir\'e fringe and its photographic negative are the same
evidence, but a dot product treats them as opposites. AQA adds a second scoring
channel based on the Born rule from quantum measurement. Each token is represented
as a unit vector, a ``pure state'', evolved by a small stack of learnable rotation
matrices (a Givens circuit), and two tokens are compared by the squared overlap
between their evolved states. This channel is mixed into the ordinary channel
through a learned, per-head, input-conditioned gate. The rotation matrices are
exactly orthogonal by construction and are parameterised by $\mathcal{O}(dL)$ angles
rather than the $\mathcal{O}(d^{2})$ weights of a normal linear layer.

\textbf{A quantum evidence basis that makes the explanation part of the
computation.} A learnable orthonormal set of basis vectors, built as the image of
the standard basis under a rotation circuit, stays exactly orthonormal at every
step of training and is measured against every spatial token. The resulting
probabilities form a token-by-mode evidence map that is used as a classification
feature \emph{and}, completely unchanged, as the explanation. The explanation is
therefore a direct read-out of the computation rather than a post-hoc approximation
of it produced by a separate model.

\textbf{Pseudo-Unseen Evidence Suppression (PUES).} If you take two training images
from two different attack families and blend them, you get something no real camera
or generator could ever produce, because no physical process makes a superposition
of a print attack and a face swap. That makes these blends a free, built-in proxy
for ``novel input''. We train the model to push its evidential output
toward zero on these blends, which teaches the uncertainty channel what ``I do not
recognise this'' should look like without ever showing the model a real unseen
attack.

We also built \textbf{evidence-grounded reasoning}. Every forward pass emits a
typed, numeric evidence record, and the natural-language explanation is generated
only from that record, either by a deterministic rule-based generator whose
hallucination rate is zero by construction, since every sentence is gated by a
numeric condition, or optionally by a hosted multimodal language model whose output
is then automatically checked against the record.

Finally, because no existing capture-based dataset marks \emph{where} in the frame
the attack evidence lives, we built a difficulty-calibrated, identity-disjoint
benchmark of eight attack families with pixel-level evidence masks, along with
identity-consistent, leave-one-attack-out and cross-domain protocols, and an
explicit calibration procedure so the benchmark cannot be trivially solved.

Section~\ref{sec:related} surveys the research threads X-UAD tries to bring
together and states the gap explicitly. Section~\ref{sec:bench} describes the
benchmark. Section~\ref{sec:method} describes the model. Section~\ref{sec:exp}
reports the experiments. Section~\ref{sec:gap} returns to the gap and explains,
piece by piece, what closes it.

% =====================================================================
\section{Related Work and the Gap}\label{sec:related}
% =====================================================================

\subsection{Physical presentation attack detection}
Early face anti-spoofing treated liveness as a texture classification problem.
Auxiliary supervision \cite{liu2018auxiliary} replaced the plain real/fake label
with a pseudo-depth map and an rPPG signal as extra training targets. Central
Difference Convolution \cite{yu2020cdcn} modified the very first layers of the
network to respond to illumination-invariant micro-texture by mixing in the local
gradient response, and it remains one of the strongest lightweight backbones for
physical attacks specifically. Pixel-wise supervision was revisited more
systematically in \cite{yu2021revisiting}. Several works then tackled the harder
problem of generalising across capture domains, meaning different cameras and
different lighting, using adversarial domain generalisation \cite{shao2019mada},
single-side domain generalisation \cite{jia2020ssdg}, and style-shuffled feature
separation \cite{wang2022ssan}. All of this work shares one assumption: that the
attack happens in front of the camera. None of it has any mechanism for a capture
that is physically genuine but was manipulated afterward.

\subsection{Digital forgery detection}
FaceForensics++ \cite{rossler2019ff} and Celeb-DF \cite{li2020celebdf} are the
benchmarks this line of work grew up around. The most transferable cue turned out
to be spectral. Generative models leave periodic artifacts from their upsampling
layers \cite{odena2016checkerboard}, which show up as distinctive peaks in the
Fourier magnitude \cite{frank2020frequency}, and this is a large part of why fairly
simple detectors generalise surprisingly well across different generator
architectures \cite{wang2020cnnspot}. The symmetric weakness applies here too. These
detectors have no concept of a re-acquisition medium and fail badly on printed or
replayed attacks.

\subsection{Unified physical-digital detection}
UniAttackData \cite{yu2024unified} is arguably the most important recent
contribution in this space. It is an identity-consistent corpus in which the same
person appears both genuine and under every attack type, so that identity itself
cannot become a shortcut, paired with a vision-language detector using
teacher-student prompting. The CVPR'24 and ICCV'25 challenges built on this
\cite{yuan2024challenge,ma2025uniattackplus}, scaling to nearly $2{,}900$ subjects
and $54$ attack types, and confirmed two things worth remembering: methods built on
large pretrained multimodal backbones dominate the leaderboard, and
single-modality CNNs get stuck above $20\%$ ACER on the harder protocol. UniDefense
\cite{zhang2024unidefense} approaches the same goal through dual-space
reconstruction. What none of these systems has is an uncertainty channel, an
evidence map, or a genuine explanation. All that comes out is a score.

\subsection{Uncertainty and open-set recognition}
Maximum softmax probability \cite{hendrycks2017msp} is a weak out-of-distribution
signal for the reason described above, since softmax is a closed-set
construction. Prior networks \cite{malinin2018prior} and evidential deep learning
\cite{sensoy2018edl}, built on subjective logic \cite{josang2016subjective}, place a
Dirichlet distribution over the simplex so that total evidence, and therefore
vacuity, becomes something a model can output. Outlier exposure
\cite{hendrycks2019oe} showed that having even a rough proxy for ``something else''
during training is worth more than any purely intrinsic uncertainty score. Open-set
single-domain generalisation has recently been explored for anti-spoofing on its own
\cite{liu2024openset}, but to our knowledge evidential uncertainty has not been
combined with unified physical-digital detection before, and no prior work builds
its outlier proxy directly out of the attack taxonomy the way PUES does.

\subsection{Explainability and MLLM-based forensics}
Attention rollout \cite{abnar2020rollout} and deletion/insertion faithfulness
testing \cite{petsiuk2018rise} are the standard tools for checking whether an
explanation reflects a model's decision. FFAA \cite{huang2024ffaa}
fine-tunes a multimodal LLM for open-world forgery analysis with a multi-answer
decision system, and M2F2-Det \cite{guo2025m2f2det} bridges CLIP and an LLM for
interpretable forged-face detection. Both produce fluent text, and both are, in a
real sense, narrating after the fact. The language model sees the image but never
the detector's internal evidence, so there is no way to check its account against
the decision it is supposedly explaining, and it is free to describe artifacts that
are not there.

\subsection{Quantum-inspired attention}
Quantum Vision Transformers \cite{cherrat2024qvt} and quantum self-attention
\cite{comajoan2024qgluon,evans2025on} replace ordinary linear projections with
parameterised quantum circuits, reporting accuracy comparable to classical vision
transformers at a small fraction of the parameters, which is essentially an
$\mathcal{O}(n^{2})\to\mathcal{O}(n)$ argument. Quantum adaptive self-attention
\cite{chen2025qasa} places a parameterised circuit inside a single encoder layer
and compares it against capacity-matched classical controls, to separate a quantum
effect from an architectural one. This line of work is motivated
primarily by parameter efficiency and is evaluated on generic image classification.
It has not been connected to a task whose evidence is literally an interference
pattern, and none of it offers an explainability story.

\subsection{Stating the gap}
Every one of the six lines above occupies a proper subset of what an operator
needs (Table~\ref{tab:gap}), and the missing pieces are not just
additive. Unifying the attack space creates the uncertainty problem, because more
families means more ways to fall outside all of them. Adding uncertainty creates the
explanation problem, because an operator asked to adjudicate a vacuous,
low-confidence sample needs to know what the model saw. And an explanation
is only worth trusting if it is tied to the same internal evidence that produced the
decision in the first place. X-UAD is designed around this dependency chain rather
than as a grab-bag union of separate modules bolted together.

\begin{table*}[t]
\centering
\caption{Capability matrix over the lines of work X-UAD unifies. \emph{Partial}
means the capability exists only for one attack super-class, or only as a proxy.}
\label{tab:gap}
\footnotesize
\setlength{\tabcolsep}{4pt}
\resizebox{\textwidth}{!}{%
\begin{tabular}{lccccccc}
\toprule
 & Physical & Digital & Second-order & Explicit & Evidence & Natural-language & Explanation bound \\
Representative work & PAs & attacks & uncertainty & unknown class & localisation & explanation & to the decision \\
\midrule
CDCN++ \cite{yu2020cdcn}, pixel-wise FAS \cite{yu2021revisiting} & yes & no & no & no & partial (depth proxy) & no & --- \\
SSAN \cite{wang2022ssan}, SSDG \cite{jia2020ssdg} & yes & no & no & no & no & no & --- \\
Frequency forensics \cite{frank2020frequency,wang2020cnnspot} & no & yes & no & no & no & no & --- \\
UniAttackDetection \cite{yu2024unified}, UniDefense \cite{zhang2024unidefense} & yes & yes & no & no & no & no & --- \\
Challenge winners \cite{ma2025uniattackplus} (CLIP/VAE, LoRA ViT-L) & yes & yes & no & no & no & no & --- \\
Evidential / open-set \cite{sensoy2018edl,liu2024openset} & partial & no & yes & yes & no & no & --- \\
FFAA \cite{huang2024ffaa}, M2F2-Det \cite{guo2025m2f2det} & no & yes & no & partial & partial & yes & no (post-hoc) \\
Quantum ViT \cite{cherrat2024qvt,evans2025on,chen2025qasa} & no & no & no & no & no & no & --- \\
\midrule
\textbf{X-UAD (ours)} & \textbf{yes} & \textbf{yes} & \textbf{yes} & \textbf{yes} & \textbf{yes} & \textbf{yes} & \textbf{yes} \\
\bottomrule
\end{tabular}}
\end{table*}

% =====================================================================
\section{The X-UAD Benchmark}\label{sec:bench}
% =====================================================================

\subsection{Why we built our own benchmark}
The capture-based corpora the unified-detection literature relies on
\cite{yu2024unified,ma2025uniattackplus} are access-restricted, and more
importantly for our purposes, none of them marks \emph{where} in the frame the
attack evidence sits. Training and scoring an evidence localiser needs
ground truth that a real capture pipeline cannot provide, because nobody
knows exactly which pixels a real printer or a real generator changed. We therefore
build a benchmark where the evidence is injected by a known, documented operator, so
its support is known exactly, by construction.

This involves a trade-off. A generated benchmark buys
pixel-level supervision, perfect identity consistency, and full reproducibility, and
it costs external validity, since an operator is a \emph{model} of a physical or
generative process rather than the process itself. We try to manage that risk in
four ways. Every operator is motivated by the actual physics or generator
architecture it stands in for. The benchmark is difficulty-calibrated so it cannot
be solved trivially. We keep a second, different capture domain purely
for cross-domain testing. And we release configuration files that map the identical
taxonomy, protocols and model code onto CASIA-FASD \cite{zhang2012casia}, OULU-NPU
\cite{boulkenafet2017oulu}, SiW-Mv2, CelebA-Spoof \cite{zhang2020celebaspoof},
FaceForensics++ \cite{rossler2019ff}, Celeb-DF-v2 \cite{li2020celebdf} and
UniAttackData \cite{yu2024unified}, so that every claim in this paper is directly
re-testable against real captured data by anyone who obtains access to those
corpora. Nothing about the model, the losses, the training loop, the
explainability code or the reasoning module is specific to our synthetic benchmark.

\subsection{Source faces and identity discipline}
\textbf{Domain A (LFW-Sim)} is built from $12{,}794$ identity-labelled, tightly
cropped face photographs of $5{,}616$ subjects from Labeled Faces in the Wild
\cite{huang2007lfw}. Splits are assigned by a stable hash of the subject identity,
so train, validation and test never share a subject. Every source face contributes
exactly one genuine sample and one sample of each of the eight attack families, so
the mutual information between identity and label is zero by construction and no
detector can cheat by memorising who a person is rather than learning what an attack
looks like. The source imagery is grayscale, and we treat that as a deliberate
constraint rather than a limitation. It removes colour-based shortcuts such as
display gamut shift, ink gamut compression and colour-transfer residue, and forces
every method, ours included, to make its decision on texture, geometry and frequency
evidence instead.

\textbf{Domain B (Wild-Sim)} applies the identical operators to $700$ in-the-wild
colour face crops at low and variable resolution with an unknown prior JPEG
compression history. It is never used for training. Its only job is to be the target
of the cross-domain protocol, and it differs from Domain A in colour, alignment,
resolution and compression history all at once.

\subsection{Attack taxonomy}
Eight classes span the two super-classes (Fig.~\ref{fig:taxonomy}). Each operator
takes a global strength parameter $s\in(0,1]$ that scales every artifact's
amplitude, and returns a soft mask $M\in[0,1]^{H\times W}$ marking which pixels
carry the injected evidence.

\begin{itemize}
\item \textbf{print} (physical): Yule--Nielsen tone transfer with dot gain, an
      ordered halftone screen at a sampled line frequency and angle, paper
      substrate grain, loss of live micro-texture, a small planar perspective
      shift, and a simulated re-acquisition chain comprising signal-dependent shot
      and read noise followed by JPEG compression.
\item \textbf{replay} (physical): display gamut and contrast compression, an RGB
      sub-pixel stripe lattice at a sampled period and tilt that aliases into a
      visible low-frequency moir\'e beat, rolling-shutter banding, specular
      backlight glare, then re-acquisition and double compression.
\item \textbf{mask3d} (physical): edge-preserving smoothing confined to the skin
      region, since silicone loses live micro-texture, with eye and mouth cut-outs
      through which the wearer's real features remain visible, a waxy albedo with
      broad specular sheen, and a seam at the mask's rim.
\item \textbf{faceswap} (digital): a donor face aligned to the target's face box,
      passed through a deliberately low generator resolution and upsampled to mimic
      the resolution mismatch of a real swap pipeline, colour-transferred, and
      alpha-blended into the inner-face region. The residual evidence is the
      blending boundary and resolution mismatch, not gross misalignment.
\item \textbf{reenact} (digital): a smooth non-rigid warp confined to the mouth and
      eye regions, with a local sharpness and contrast inconsistency inside the
      manipulated area only.
\item \textbf{synth} (digital): the classic generator signature, namely
      over-smoothed skin combined with over-sharpened detail, a
      transposed-convolution checkerboard at a sampled stride
      \cite{odena2016checkerboard}, and an explicit azimuthally modulated radial
      peak injected into the Fourier magnitude \cite{frank2020frequency}.
\item \textbf{adversarial} (digital): not random noise, but a genuine
      $\ell_\infty$ PGD perturbation ($\varepsilon=6/255$, $7$ steps)
      \cite{madry2018pgd} crafted against an independently trained surrogate
      detector, a small VGG-style network that shares no weights, architecture or
      random seed with any model we evaluate, and is fit only on the training
      split. The perturbations seen at validation and test time are therefore a
      genuine transfer, that is, black-box, attack. The surrogate itself is flipped
      to ``genuine'' on $38.8\%$ of the crafted samples. Ground truth for
      localisation is the smoothed, normalised magnitude of the perturbation itself.
\end{itemize}

\begin{figure}[t]
\centering
\includegraphics[width=\linewidth]{figures/taxonomy.png}
\caption{The unified taxonomy. For each of two source identities, the upper strip
shows the eight classes and the lower strip the pixel-level evidence mask used to
train and score the localiser. Physical operators (print, replay) modify the whole
capture; region-level operators (3-D mask, face-swap, reenactment, adversarial)
have a localised support. The region-level IoU of Table~\ref{tab:loc} is
averaged over 3-D mask, face-swap and reenactment.}
\label{fig:taxonomy}
\end{figure}

\subsection{Difficulty calibration}
A simulated benchmark is worthless if a simple linear probe can solve it outright.
We fix each operator's strength through a single documented procedure: sweep $s$ and
pick the largest value at which a reference detector (CDCN++, our strongest
lightweight texture baseline) does \emph{not} saturate, leaving measurable headroom
on both the binary and the eight-way task. This lands on $s=0.62$, where the
reference detector reaches $1.92\%$ ACER but only $91.0\%$ eight-way accuracy and
$3.74\%$ expected calibration error, with a worst-case per-family APCER of $2.67\%$.
The binary problem therefore sits close to, but not at, the ceiling, while
family discrimination, calibration and unseen-attack behaviour all remain wide open.
This mirrors, on a benchmark we fully control, the pattern reported for
UniAttackData \cite{yu2024unified}: a saturated seen-attack protocol next to a
wide-open unseen-attack one.

\subsection{Protocols}
\textbf{P1 (seen attacks).} All eight families appear in train, validation and test,
with identity-disjoint splits. This measures closed-set unified detection and family
discrimination.

\textbf{P2 (leave-one-attack-out).} One family is withheld entirely during training
and appears only at test time. We report ACER on the full test set, the detection
rate of the held-out family at the fixed operating threshold, and the \emph{novelty
AUROC}: the AUROC with which the model's own uncertainty channel ranks the held-out
family \emph{above the attack families it was trained on}, which is to say whether
the model knows that this is a kind of attack it has not seen.

\textbf{P3 (cross-domain).} Train on Domain A, test unchanged on Domain B, with zero
adaptation.

All thresholds are fixed at the equal-error point of the validation split and
applied unchanged to test. Every metric follows ISO/IEC 30107-3 \cite{iso30107}:
APCER is the worst case across attack families, BPCER is the genuine-rejection rate,
and ACER is their average.

% =====================================================================
\section{Method}\label{sec:method}
% =====================================================================

\subsection{Overview}
X-UAD maps a face crop $\bs{x}\in\R^{3\times H\times W}$ to a tuple
\[
\bs{x}\;\longmapsto\;\bigl(\bs{\alpha},\;u,\;\bs{M},\;\bs{\Pi},\;r,\;\tau\bigr),
\]
where $\bs{\alpha}\in\R^{K}_{>0}$ are Dirichlet concentrations over the $K=8$
classes, $u\in[0,1]$ is the vacuity (a single scalar answering ``how much do I have
to go on''), $\bs{M}$ is the spoof-evidence map, $\bs{\Pi}$ is the quantum evidence
spectrum over spatial tokens and basis modes, $r$ is the risk score, and $\tau$ is
the natural-language rationale. All six come out of a single forward pass
(Fig.~\ref{fig:arch}). A central-difference stem and a spectral branch produce
tokens; a stack of Adaptive Quantum Attention blocks processes them; the quantum
evidence basis is measured against the output tokens; an evidential head, a
localisation decoder and a risk calibrator read everything out; and the reasoning
module turns the record into text.

\begin{figure*}[t]
\centering
\includegraphics[width=0.94\textwidth]{figures/architecture.png}
\caption{X-UAD. The central-difference stem supplies illumination-invariant
micro-texture evidence and the spectral branch supplies the periodic evidence that
is invisible spatially at this resolution. Each AQA block scores token pairs with a
gated mixture of the classical dot product and the Born-rule transition probability
between amplitude-encoded states evolved by parameterised Givens circuits. The
learnable orthonormal evidence basis is measured against the output tokens; its Born
probabilities feed the classifier \emph{and} are exported unchanged as the
explanation, which is what binds the rationale to the decision.}
\label{fig:arch}
\end{figure*}

\subsection{Central-difference and spectral evidence stem}
Central Difference Convolution \cite{yu2020cdcn} interpolates between an ordinary
convolution and its local gradient response,
\begin{equation}
y(p_{0}) \;=\; \sum_{p\in\mathcal{N}} w(p)\,x(p_{0}+p)
\;-\;\theta \Bigl(\sum_{p\in\mathcal{N}} w(p)\Bigr) x(p_{0}),
\label{eq:cdc}
\end{equation}
which amplifies the intensity-invariant part of the response. Note that the extra
term is just a single $1\times1$ convolution, not a second full convolution, so it
is nearly free. Four strided CDC stages take a $112^{2}$ input down to $7^{2}$,
exposing the intermediate $28^{2}$ and $14^{2}$ maps to the localisation decoder
later on.

Periodic structure, meaning the display lattice behind a replay attack or the
upsampling checkerboard behind a generated face, is not represented well in a
$7\times7$ spatial token grid; it gets aliased away. We add a compact spectral
branch. The log-magnitude DFT of the luminance channel is pooled onto a polar grid
($16$ radii $\times$ $12$ angles), which is translation-invariant and
rotation-equivariant by construction, and a small two-layer MLP encodes this profile
into four spectral tokens that get appended to the spatial tokens. The same pooled
profile is also exported into the evidence record, so the radial peak frequency, its
prominence, and the azimuthal anisotropy become concrete, quotable numbers in the
eventual explanation.

\subsection{Adaptive Quantum Attention}
\textbf{Amplitude encoding.} For a head with dimension $d$, queries and keys are
$\ell_{2}$-normalised to unit vectors,
$\ket{\psi^{Q}_{i}} = \bs{q}_{i}/\lVert\bs{q}_{i}\rVert$ and
$\ket{\psi^{K}_{j}} = \bs{k}_{j}/\lVert\bs{k}_{j}\rVert$, a real-valued amplitude
encoding of $d=2^{n}$ amplitudes across $n$ notional qubits.

\textbf{Parameterised rotation circuit.} A depth-$L$ circuit acts on a ring of
amplitudes. Layer $\ell$ rotates disjoint adjacent pairs, alternating which pairs it
touches between layers, each rotation governed by a single learnable angle:
\begin{equation}
G(\vartheta)=\begin{pmatrix}\cos\vartheta & -\sin\vartheta\\
\sin\vartheta & \cos\vartheta\end{pmatrix},
\qquad
U(\bs{\vartheta})=\prod_{\ell=1}^{L}\ \bigotimes_{k} G(\vartheta_{\ell k}).
\label{eq:givens}
\end{equation}
This is the real-amplitude analogue of a hardware-efficient ansatz built from
nearest-neighbour two-qubit gates with periodic boundary conditions. Alternating the
two offsets keeps the circuit connected, so the set of reachable rotations grows
toward the full orthogonal group as $L$ increases. Three things matter about this
construction. First, it is exactly unitary, since $U$ is a product of rotations, so
it stays orthogonal to machine precision at every step of training. We measured
$\lVert\,\lVert U\bs{x}\rVert - \lVert\bs{x}\rVert\,\rVert_{\infty} < 10^{-6}$
throughout, so no re-orthogonalisation or penalty term is ever needed. Second, it is
parameter-efficient, costing $L\lfloor d/2\rfloor$ angles, that is
$\mathcal{O}(dL)$, rather than the $\mathcal{O}(d^{2})$ of a dense projection
matrix, the same asymptotic argument made for quantum self-attention elsewhere
\cite{cherrat2024qvt,evans2025on}. Third, because it exactly preserves the vector's
norm, the resulting similarity score below is a genuine probability rather than just
a bounded number.

\textbf{The Born channel.} Queries and keys are evolved by two \emph{different}
circuits. This matters, because if they shared one circuit the rotation would cancel
out of the inner product entirely and add nothing. The similarity score is the
Born-rule transition probability,
\begin{equation}
S_{ij} \;=\; \bigl\lvert \bra{\psi^{Q}_{i}} U_{Q}^{\!\top} U_{K} \ket{\psi^{K}_{j}} \bigr\rvert^{2}
\;\in\;[0,1].
\label{eq:born}
\end{equation}
This is a quadratic similarity. It is invariant to a sign flip of either state, and
it is maximised equally by alignment \emph{or} anti-alignment. That is the
invariance an interference pattern has. A moir\'e fringe and its negative carry the
same forensic information, and it is precisely the invariance an ordinary dot
product lacks.

\textbf{Adaptive gating.} Not every kind of evidence is interference-like. A
blending boundary, for instance, is really just a plain edge. A per-head gate,
conditioned on simple statistics of the token set, decides the mixture between the
two channels:
\begin{align}
\bs{g} &= \sigma\bigl(\bs{W}_{g}\,[\,\mathrm{mean}_t(\bs{X}) \,\Vert\, \mathrm{std}_t(\bs{X})\,] + \bs{b}_{g}\bigr) \in [0,1]^{H_{\text{heads}}},\\
\bs{A} &= \mathrm{softmax}\Bigl( (1-g_{h})\tfrac{\bs{Q}\bs{K}^{\!\top}}{\sqrt{d}} \;+\; g_{h}\,\beta\, \bs{S} \Bigr),
\label{eq:gate}
\end{align}
with $\beta$ a learnable temperature. $\bs{W}_{g}$ is initialised to zero and
$\bs{b}_{g}$ to $-2$, so the gate starts at $g\approx0.12$, and every block is
wrapped in LayerScale \cite{touvron2021cait} with residual branches multiplied by a
learnable per-channel gain initialised to $10^{-2}$. In practice this means the
network starts life as basically its own convolutional sub-network, the
central-difference stem plus a plain classifier, and only gradually recruits
attention, with the Born channel inside attention only earning gradient through its
residual as training proceeds. This matters more than it might sound: in controlled
probes at matched budget, removing LayerScale let the freshly initialised attention
inject noise into the stem and slowed convergence relative to either component
trained on its own. The converged gate value $g$ is reported directly to the
operator, which makes the mechanism itself auditable rather than a black box.

\textbf{Measurement-entropy regularisation.} Let
$\mathcal{H}=\mathbb{E}_{t}\bigl[-\sum_{j}A_{ij}\log A_{ij}\bigr]$ be the mean row
entropy of the attention matrix. Too low means attention has collapsed onto a single
token and stopped integrating any real evidence; too high means it has diffused into
a uniform mixture and is carrying nothing at all. A parameter-free two-sided hinge
keeps it inside a band:
\begin{equation}
\mathcal{L}_{\text{coh}} = \bigl[\lambda_{\text{lo}}\log T - \mathcal{H}\bigr]_{+}
+ \bigl[\mathcal{H} - \lambda_{\text{hi}}\log T\bigr]_{+},
\end{equation}
with $(\lambda_{\text{lo}},\lambda_{\text{hi}})=(0.30,0.80)$.

\subsection{Quantum evidence basis}
Let $\{\ket{e_{b}}\}_{b=1}^{B}$ be the first $B$ computational basis states and
$V(\bs{\phi})$ a second rotation circuit of the same family described above. The
evidence basis is $\ket{\phi_{b}} = V(\bs{\phi})\ket{e_{b}}$, which stays exactly
orthonormal for every value of $\bs{\phi}$. Measuring the output token states
against it gives Born probabilities
\begin{equation}
\Pi_{tb} \;=\; \bigl\lvert \braket{\phi_{b}}{\psi_{t}} \bigr\rvert^{2},
\qquad \textstyle\sum_{b}\Pi_{tb}\le 1 .
\label{eq:basis}
\end{equation}
$\bs{\Pi}$ does two jobs at once, and that dual role is really the whole design
point. Its token-wise average is concatenated straight into the classification
feature vector, and the same tensor, not an approximation of it and not some
separate gradient-based surrogate map, is exported unchanged as the explanation. A
lightweight auxiliary classifier trained on $\overline{\bs{\Pi}}$ with ordinary
cross-entropy pushes each mode to become class-selective, so that after training
every mode can be \emph{named} after the family it responds to most
(Table~\ref{tab:modes}). Those names are read off the trained model rather than
assigned by hand.

\subsection{Decoupled evidential head and vacuity}
A Dirichlet distribution is completely described by a direction on the simplex and
a concentration. The standard evidential formulation \cite{sensoy2018edl} predicts
all $K$ concentrations jointly; we predict the two factors separately, which turned
out to matter quite a bit in practice. The direction (which class) and magnitude
(how much evidence) are
\begin{align}
\bs{\pi} &= \mathrm{softmax}(\bs{W}_{c}\,\bs{h}), \label{eq:dir}\\
E &= \exp\bigl(\mathrm{clip}(f_{e}(\mathrm{sg}[\bs{z}]))\bigr), \label{eq:mag}
\end{align}
from which the Dirichlet concentrations, total evidence, and vacuity follow:
\begin{equation}
\bs{\alpha} = \bs{1} + E\,\mathrm{sg}[\bs{\pi}], \qquad S = K + E, \qquad
u = \frac{K}{S},
\label{eq:vacuity}
\end{equation}
where $\mathrm{sg}[\cdot]$ denotes a stop-gradient. The subjective-logic reading
stays intact: belief masses are $b_{k}=E\pi_{k}/S$, vacuity is $u=K/(K+E)$, and the
projected probability under a uniform base rate is
$\hat{p}_{k}=(1+E\pi_{k})/(K+E)$. We report $\bs{\pi}$ as the belief direction and
$u$ as the uncertainty, the same separation subjective logic draws between what I
believe and how much I have to go on.

We verified three consequences of this decoupling empirically, in budget-matched
probes. First, the class decision keeps a clean log-loss gradient, since
$\arg\max_{k}\hat{p}_{k}=\arg\max_{k}\pi_{k}$, so training the class decision
reduces to ordinary cross-entropy on $\bs{W}_{c}\bs{h}$. The joint parameterisation,
by contrast, routes every gradient through $\bs{\alpha}$, whose Brier fit saturates
early and whose usual KL regulariser drives the pre-activation negative until the
softplus derivative collapses. We call this \emph{evidence starvation}, and we
measured it directly: swapping the joint head for the decoupled one recovered
several points of eight-way accuracy at an identical training budget. Second, adding
uncertainty cannot cost detection accuracy, because the magnitude branch only ever
reads a stop-gradient copy of the shared representation. Letting its own objective
(especially PUES, below) reshape those features would be circular, and we measured
it to be costly when we tried it. Third, suppressing evidence becomes a strictly
one-dimensional problem: because all of the concentration lives in the single scalar
$E$, an out-of-taxonomy constraint cannot distort the class direction at all.

Training combines a cross-entropy term on the direction, a Brier-style fit on the
concentration, and a term that ties the magnitude to the model's own correctness:
\begin{equation}
\mathcal{L}_{\text{ev}} = \mathrm{CE}(\bs{\pi},y)
\;+\; \tfrac{1}{2}\,\mathbb{E}\bigl[\lVert \bs{y}-\hat{\bs{p}}\rVert^{2}\bigr]
\;+\; \lambda_{c}\,\mathrm{BCE}\bigl(1-u,\;\mathbb{1}[\hat{y}=y]\bigr),
\label{eq:evloss}
\end{equation}
with $\lambda_{c}=0.4$. We use the expected sum-of-squares (Brier) form of the fit
rather than the more standard digamma-based maximum likelihood, because the latter's
gradient effectively vanishes while the Dirichlet is still near-uniform, which is
the regime a freshly initialised model starts in.

Two details here were forced on us by direct measurement rather than chosen for
elegance, and we think they matter enough to spell out. We drop the usual KL
regulariser entirely. In the joint parameterisation the KL term is what stops
evidence from being placed on the wrong classes, but here the wrong-class evidence
is already $E\pi_{k}$, which cross-entropy drives to zero on its own by sharpening
$\bs{\pi}$. The KL term therefore has nothing left to act on except the shared
magnitude $E$, and leaving it in causes short-schedule models to saturate at
$u\approx1$ everywhere and the open-set channel to die. We also added the
correctness term specifically because it is what makes vacuity useful
\emph{quickly}. The Brier fit does reward concentration, but only indirectly, so on
a short schedule the magnitude head barely moves from its initialisation without it.
Training $1-u$ directly as ``the probability this decision is correct'' makes the
channel usable within about two epochs. In a controlled leave-one-attack-out probe,
adding this term moved the AUROC with which vacuity separates a held-out family from
genuine samples from $25.6\%$ (effectively inverted) to $73.9\%$.

\subsection{Pseudo-Unseen Evidence Suppression}
An intrinsic uncertainty score is known to be weaker than one trained against an
actual proxy for ``something else'' \cite{hendrycks2019oe}, but for attack detection
there is no natural pool of outliers lying around. We build one out of the taxonomy
itself. Take two training images $\bs{x}_{a},\bs{x}_{b}$ from \emph{different}
attack families and mix them with $\lambda\sim\mathrm{Beta}(2,2)$:
$\tilde{\bs{x}}=\lambda\bs{x}_{a}+(1-\lambda)\bs{x}_{b}$. This blend carries a
superposition of two forensic signatures that no single physical or generative
process could ever produce, so it is outside the taxonomy by construction. We
constrain its evidence to sit well below that of real samples:
\begin{equation}
\mathcal{L}_{\text{PUES}} =
\frac{\bigl[\;\overline{\lVert\tilde{\bs{e}}\rVert_{1}}
      \;-\;\rho\,\overline{\lVert\bs{e}\rVert_{1}}^{\,\text{sg}}\;\bigr]_{+}}
     {\overline{\lVert\bs{e}\rVert_{1}}^{\,\text{sg}}+1},
\qquad \rho = 0.25,
\end{equation}
where bars denote batch means. The \emph{relative} form here matters and was arrived
at empirically. An absolute target ($\lVert\tilde{\bs{e}}\rVert_{1}\to0$) is
trivially satisfied by the degenerate solution of emitting zero evidence
\emph{everywhere}, collapsing the classifier to a uniform guess. We saw exactly that
failure, $12.5\%$ accuracy over eight classes, in an early version. Tying the
constraint to the batch's own in-taxonomy evidence removes that shortcut entirely,
because the fit term still demands sharply peaked posteriors on the real samples in
the same batch. This is worth distinguishing from mixup \cite{zhang2018mixup}, which
asks for a mixed \emph{label}. PUES asks for the \emph{absence of belief}, which
only makes sense because the head is evidential in the first place.

\subsection{Evidence localisation}
A lightweight feature-pyramid decoder fuses the $7^{2}$ token map with the $14^{2}$
and $28^{2}$ intermediate stem maps and predicts $\bs{M}$ at half input resolution,
supervised with a BCE plus soft-Dice objective against the operator's own
ground-truth mask. Unlike depth-map supervision
\cite{liu2018auxiliary,yu2021revisiting}, this target is the actual support of the
forensic evidence, and it is well-defined for digital attacks just as much as for
physical ones.

\subsection{Risk calibration}
Attack families are not equally harmful. With per-family severity weights
$\bs{\varsigma}$, an unknown-attack severity $\varsigma_{\text{unk}}$, and evidence
area $a=\overline{M}$,
\begin{equation}
r = (1-u)\sum_{k}\hat{p}_{k}\varsigma_{k} \;+\; \lambda_{u}\,u\,\varsigma_{\text{unk}}
\;+\; \lambda_{a}\,a .
\label{eq:risk}
\end{equation}
The first term is the expected harm under the posterior, damped by how much the
model knows. The second charges an explicit unknown-attack premium, so a
confidently-unknown sample gets escalated rather than quietly waved through, which is
the behaviour an operator needs and a plain softmax pipeline structurally cannot
give. The third term rewards spatially corroborated evidence. All coefficients, and
the band edges that turn $r$ into a \texttt{Low/Medium/High/Critical} label, are
treated as deployment policy rather than learned parameters. No training signal
touches them, deliberately, since risk is the one part of the pipeline an operator
has to be able to audit and adjust, and it is the last place we want to hide a
silently-fitted parameter. A vacuity of $u\ge0.55$ overrides the verdict and refers
the sample to a human reviewer.

\subsection{Evidence-grounded reasoning}
Every forward pass emits a typed \textbf{evidence record}: the posterior
$\hat{\bs{p}}$, the vacuity $u$, the risk score and band, the top facial regions of
$\bs{M}$ and their mean response, the area and spatial dispersion of $\bs{M}$, the
three strongest quantum modes with their learned names and activations, the spectral
descriptors, and the mean adaptive gate value. Language is generated \emph{only}
from this record, through either of two interchangeable back-ends. The
\emph{grounded realiser}, used for every result in this paper, is a deterministic
generator where every sentence is emitted by a rule whose guard is a numeric
condition on the record, so every sentence traces back to a measured quantity, the
hallucination rate is zero by construction, and it costs under a millisecond with no
network call at all. The \emph{constrained-MLLM back-end} instead hands the record to
a hosted multimodal language model with an instruction to paraphrase only, never to
add anything, and then runs an automatic verification pass that re-checks every
numeric claim in the generated text against the record and flags anything
unsupported. This inverts the usual dependency in MLLM-based forensics
\cite{huang2024ffaa,guo2025m2f2det}, where the language model sees the image but not
the detector's internal state. Here the detector's state \emph{is} the input, so the
explanation cannot contradict the decision it is describing. A verbatim example is
shown in Fig.~\ref{fig:report}.

\subsection{Training objective}
\begin{equation}
\mathcal{L} = \mathcal{L}_{\text{ev}} + w_{m}\mathcal{L}_{\text{loc}}
+ w_{b}\mathcal{L}_{\text{basis}} + w_{s}\mathcal{L}_{\text{sup}}
+ w_{p}\mathcal{L}_{\text{PUES}} + w_{c}\mathcal{L}_{\text{coh}},
\end{equation}
with $(w_{m},w_{b},w_{s},w_{p},w_{c})=(0.6,0.25,0.3,0.15,0.02)$ throughout.
$\mathcal{L}_{\text{sup}}$ is a cross-entropy on the marginal over the three
super-classes (genuine, physical, digital), obtained by summing $\hat{p}$ within each
group. It imposes the taxonomy's own hierarchy, so a mix-up between two families
inside the same super-class is penalised less harshly than one that crosses the
physical/digital boundary.

For the main seen-attack results we add three terms that act on the predicted
distribution only and add no parameters. The classification term is modulated by
a focal factor $(1-\hat{p}_{y})^{\gamma}$ with $\gamma=2$. A binary cross-entropy
$\mathcal{L}_{\text{bin}}$ on the attack probability $1-\hat{p}_{0}$ (weight $0.5$)
keeps the attack/genuine decision, which ACER is computed from, from being
traded away for family accuracy. A pair penalty
\begin{equation}
\mathcal{L}_{\text{pair}} = \sum_{(c,c')} w_{cc'}\,
\mathbb{E}_{i:\,y_i=c}\bigl[\hat{p}_{c'}(x_i)\bigr]
\end{equation}
discourages the specific confusions seen in development: face-swap and
reenactment, and 3-D mask and synthetic, in both directions with $w=0.15$; replay
and print, and face-swap and print, with $w=0.05$. Each pair is averaged over its
own samples, so adding a pair does not dilute the others. The baselines get the
same focal factor and binary term, and a pooled version of
$\mathcal{L}_{\text{pair}}$ over the first two pairs with a single weight of $0.15$
(Section~\ref{sec:exp}).

% =====================================================================
\section{Experiments}\label{sec:exp}
% =====================================================================

\subsection{Setup}
We train every model ourselves, from scratch and without pretrained weights, so
the comparison is between architectures and objectives rather than pretraining
corpora. The training budget differs between protocols, and models are only
compared within a table.

For P1 (Table~\ref{tab:main}) X-UAD, CDCN++ and ResNet-18 are trained for $12$
epochs on the full $19{,}200$-sample training split with the recipe of
Section~\ref{sec:method}: focal factor, binary term and pair penalty. The
checkpoint kept is the epoch with the best validation eight-way accuracy plus
macro-F1 among the epochs whose validation ACER is within $0.5$ points of the best
seen so far. For X-UAD we trained seven variants under this budget (with and
without the binary term, three forms of the pair penalty, a hierarchical head, two
region-pooled read-outs) and report the one with
the lowest validation ACER. The baselines were each trained once, so X-UAD
received more tuning, and the reader should weigh its P1 numbers with that in
mind. ViT-Tiny was not retrained under this recipe: at the original $10$-epoch
budget it stays near chance, and its row in Table~\ref{tab:main} is that run.

The other protocols keep the original recipe, without the three extra terms. P3
reuses the $10$-epoch P1 checkpoints. P2 and its open-set ablation use $5$ epochs
on a fixed $4{,}800$-sample subset ($4{,}382$ once the held-out family is removed),
which keeps the sixteen leave-one-attack-out runs affordable on CPU, and the
component ablation uses $5$ epochs on $3{,}592$ samples.

All runs use batch size $48$, AdamW ($\text{lr}=2.5\times10^{-3}$, weight decay
$0.02$) with a one-cycle schedule, gradient clipping at $3.0$, and $112\times112$
random crops from the $128$\,px store with horizontal flips and mild photometric
jitter. There are no geometric warps, since they would destroy the periodic
signatures we want the models to find. The operating threshold is the equal-error
point of the validation split, applied unchanged to test. Everything runs on CPU.
All results are from a single seed (0). Every model is scored on the same test
set, so we compare models with McNemar's exact test on paired errors, but that
says nothing about variance across seeds.

The baselines are \textbf{ResNet-18} \cite{he2016resnet}, a width-matched residual
CNN and the classical binary detector; \textbf{ViT-Tiny} \cite{dosovitskiy2021vit},
a patch-16 transformer with ordinary attention, included as the plain-transformer
control from the same architecture family as the ViT-B/16 entry that underperforms
in \cite{yu2024unified}; and \textbf{CDCN++} \cite{yu2020cdcn,yu2021revisiting}, a
central-difference trunk with pixel-wise map supervision. CDCN++ is the strongest
lightweight texture detector available to us and the one that shares our stem, so
comparing against it isolates what the transformer, the Born channel, the
evidential head and the reasoning stack add.

\subsection{Seen-attack unified detection (P1)}

\begin{table*}[t]
\centering
\caption{Protocol P1 (seen attacks, identity-disjoint, ID-consistent), $12$-epoch
recipe, seed $0$, $9{,}600$ test samples. ACC$_8$ is eight-way family accuracy;
ECE$_{\text{TS}}$ is ECE after temperature scaling fitted on the validation split;
the other metrics follow ISO/IEC 30107-3 \cite{iso30107}. Bold marks the best of
the three detectors trained under the common recipe.
$^{\dagger}$ViT-Tiny: original $10$-epoch recipe, not retrained; its low ECE reflects
a detector operating near chance.}
\label{tab:main}
\footnotesize
\input{tables/main}
\end{table*}

P1 is close to saturated for every model with a convolutional front end
(Table~\ref{tab:main}). ResNet-18, CDCN++ and X-UAD all reach at least $99.95\%$
AUC and at most $1.9\%$ ACER, which matches what the unified-detection literature
reports on its own seen-attack protocol \cite{yu2024unified}. We tested each gap
with McNemar's exact test on the paired errors over the same $9{,}600$ samples.

On the attack/genuine decision X-UAD has the lowest ACER, $0.42\%$ against $0.46\%$
for CDCN++, but the two models disagree on only $22$ test samples, $10$ in X-UAD's
favour and $12$ against ($p=0.83$). We read this as a tie. Against ResNet-18 the
difference is clear: $74$ discordant samples in X-UAD's favour against $12$
($p=4.7\times10^{-12}$). ResNet-18 misses $3.1\%$ of face-swaps at the validation
threshold (Table~\ref{tab:pertype}); X-UAD's worst family is replay, at $0.33\%$.

The eight-way family decision goes the other way. ResNet-18 reaches $96.93\%$ and
CDCN++ $95.19\%$, against $94.06\%$ for X-UAD, and both gaps are significant ($443$
against $168$ discordant samples, $p=1.6\times10^{-29}$; $349$ against $241$,
$p=1.0\times10^{-5}$). Fig.~\ref{fig:conf} shows where X-UAD's family errors go:
the largest are synthetic predicted as 3-D mask ($7.5\%$),
face-swap and reenactment confused in both directions ($5.6\%$ and $5.4\%$), and
3-D mask predicted as synthetic ($4.6\%$), which are the pairs the pair penalty
targets. Fig.~\ref{fig:roc} gives the ROC and DET curves.

We do not find a calibration advantage for X-UAD's class posterior. Before any
correction ResNet-18 has the lowest ECE ($2.35\%$, against $5.07\%$ for CDCN++ and
$5.17\%$ for X-UAD). After temperature scaling on the validation split the three
reach $0.34\%$, $0.46\%$ and $0.60\%$ respectively, the same order as before. What X-UAD adds here is a different
quantity, the Dirichlet vacuity, which neither baseline has and which we evaluate
as an open-set score in Section~\ref{sec:p2}.

The plain patch-16 transformer fails outright, at $49.8\%$ ACER. With
$16\times16$ patches it never sees the high-frequency texture that the physical
families live in, the same weakness that puts ViT-B/16 behind ResNet-50 in
\cite{yu2024unified}, here in a starker form. It is the reason X-UAD puts a
central-difference stem under its transformer rather than tokenising raw pixels.

\begin{table}[t]
\centering
\caption{Per-family APCER (\%) at the validation-fixed threshold, protocol P1. The
ISO APCER is the \emph{worst} column.}
\label{tab:pertype}
\small
\resizebox{\linewidth}{!}{\input{tables/per_type}}
\end{table}

\begin{figure}[t]
\centering
\includegraphics[width=\linewidth]{figures/roc_det.png}
\caption{ROC in the low-false-alarm regime that matters operationally, and the DET
curve, protocol P1, $12$-epoch recipe.}
\label{fig:roc}
\end{figure}

\begin{figure}[t]
\centering
\includegraphics[width=0.82\linewidth]{figures/confusion.png}
\caption{X-UAD family confusion (row-normalised, \%), protocol P1, $12$-epoch
recipe. The largest errors are synthetic$\to$3-D mask ($7.5\%$) and
face-swap$\leftrightarrow$reenactment ($5.6\%$, $5.4\%$). Two of the
larger errors cross the physical/digital boundary: synthetic$\to$3-D mask and
face-swap$\to$print ($4.0\%$).}
\label{fig:conf}
\end{figure}

\subsection{Unseen attack families (P2)}\label{sec:p2}
This is the operationally decisive question. What happens when an attack family the
model has never seen shows up at test time? Table~\ref{tab:loao} holds out
one family at a time and reports four separate things: the ISO error on the full test
set, the eight-way family accuracy, the detection rate of the held-out family at the
unchanged threshold, and the novelty AUROC, the AUROC with which the model's
uncertainty ranks the held-out family above the attack families it was trained on.

\begin{table}[t]
\centering
\caption{Protocol P2, leave-one-attack-out. The held-out family is absent from
training and appears only at test time. Means are over all four folds. Novelty
AUROC ranks the held-out family against the \emph{seen attack} families using the
model's uncertainty score; values below $50$ indicate an inverted ranking.}
\label{tab:loao}
\small
\resizebox{\linewidth}{!}{\input{tables/loao}}
\end{table}

Every method's mean ACER is in the $29$ to $41\%$ range here, against under $2\%$
in P1. That gap should \emph{not} be read as the cost of withholding a family
alone: P2 is trained on a quarter of the data for half the epochs, and the
component ablation shows the full X-UAD already reaching $38\%$ ACER at a similar
reduced budget with \emph{every} family seen (Table~\ref{tab:abl}). How much of the
P2 error is due to the held-out family and how much to the budget is not separated
by our experiments. The comparisons that are valid are between models within
Table~\ref{tab:loao}, which all share the same budget. Within that table, X-UAD
rejects the unseen family most often on average ($77.5\%$ against $74.7\%$ and
$70.2\%$, Fig.~\ref{fig:loao}), and its uncertainty is a purpose-built channel
rather than an incidental by-product of a softmax. Its novelty AUROC clearly leads on
two of the four folds, replay ($71.0$ against $49.3$ and $35.3$) and face-swap
($87.9$ against $81.7$ and $74.4$), and trails on the other two, landing level with
CDCN++ on the mean ($63.09$ against $63.15$) and well ahead of ResNet-18.

We report the ``is it rejected'' question and the ``does uncertainty rank it above
known attacks'' question separately and deliberately, because they answer different
questions for an operator. Whether the held-out family is rejected at all is the
operational question; whether uncertainty ranks it above the seen attacks is the
analyst's question, the one that flags that a new attack vector has
appeared. The adversarial fold makes the distinction concrete. X-UAD rejects
$98.2\%$ of a family it never trained on while still assigning it low vacuity,
because the perturbation sits on top of an attack whose underlying evidence the
model does recognise. That fold is also where our uncertainty channel is weakest, at
$43.0$ novelty AUROC, meaning the ranking is mildly inverted. We have no satisfying
account of why.

The comparison above is across models. A cleaner test of whether the evidential
channel itself adds anything is to score the \emph{same} X-UAD models with their
own maximum-softmax confidence instead of vacuity. Averaged over the four folds,
vacuity gives $63.1$ novelty AUROC against $55.9$ for X-UAD's own softmax score.
Almost all of that gain comes from the replay fold ($71.0$ against $39.5$); it is
small on face-swap ($+1.2$) and adversarial ($+2.3$), and negative on 3-D mask
($-6.4$). The per-group means explain the two failing folds. On 3-D mask the
held-out family gets the same vacuity as the seen attacks ($0.31$ against
$0.32$), so vacuity does not see it. On adversarial it gets \emph{less}
($0.39$ against $0.43$), which is the inversion.

\begin{figure}[t]
\centering
\includegraphics[width=\linewidth]{figures/loao.png}
\caption{Leave-one-attack-out. Left: ISO ACER of each fold on the full test set,
held-out family included. Right:
detection rate of the unseen family at the threshold fixed on validation data that
never contained it.}
\label{fig:loao}
\end{figure}

\subsection{Cross-domain transfer (P3)}
Domain B differs from Domain A in colour, alignment,
resolution and compression history all at once, and every single method collapses
under the shift (Table~\ref{tab:cross}). With the threshold fixed on the source
domain, BPCER runs to $89$--$100\%$, meaning the detectors reject almost every genuine
face from the new domain outright. Ranked by threshold-free AUC, CDCN++ retains the
most signal at $65.96$, while ResNet-18 ($46.78$), ViT-Tiny ($50.18$) and X-UAD
($50.90$) all sit at or below chance. The adversarial class is excluded here: it was
crafted from Domain-A images against a Domain-A surrogate, so it is not cross-domain
data.

\begin{table}[t]
\centering
\caption{Protocol P3, cross-domain transfer A\,$\rightarrow$\,B, no adaptation.
Every method degrades heavily.}
\label{tab:cross}
\footnotesize
\resizebox{\linewidth}{!}{\input{tables/crossdomain}}
\end{table}

We draw two conclusions. First, this protocol is harsher than the cross-dataset protocols
usually reported elsewhere, because the shift here is joint across four factors
rather than a single one. Second, none of the mechanisms
in this paper, ours very much included, addresses capture-domain shift. That is a
different problem from attack-family shift, and it would need the
domain-generalisation machinery of \cite{wang2022ssan,shao2019mada,jia2020ssdg}
layered on top. X-UAD does not lead this table and we make no claim that it does.

\subsection{Calibration and evidence localisation}
Fig.~\ref{fig:calib}(a) shows the reliability diagrams behind the ECE column of
Table~\ref{tab:main}. All three detectors are somewhat miscalibrated before
correction, and temperature scaling \cite{guo2017calibration} removes most of it for
each of them, so calibration of the class posterior does not separate the models.
Panel (b) shows the uncertainty each model can offer an operator. For the softmax
baselines that is $1-\max_k p_k$, a by-product of a closed-set distribution. For
X-UAD it is the vacuity $u=K/S$, which is trained as an output in its own right and,
as Table~\ref{tab:loao} shows, ranks unseen families above the seen attack families
on two of the four folds.

\begin{figure}[t]
\centering
\includegraphics[width=\linewidth]{figures/calibration.png}
\caption{Protocol P1, $12$-epoch recipe. (a) Reliability before temperature
scaling. (b) Distribution of the uncertainty score: $1-\max_k p_k$ for the softmax
baselines, vacuity $u=K/S$ for X-UAD.}
\label{fig:calib}
\end{figure}

Table~\ref{tab:loc} reports intersection-over-union against the generator's own
mask. We separate out the families whose evidence occupies a sub-region of
the frame (3D mask, face-swap, reenactment) from the ones that alter the whole
capture (print, replay, synthetic), since the latter have near-constant masks and a
trivially high IoU that would tell us nothing about localisation quality. The
region-level mean is the number that measures whether the localiser works. Both
models are scored on all $9{,}600$ test samples. X-UAD reaches $80.1\%$ region-level IoU against $73.2\%$ for
CDCN++, and is ahead on each of the three region-level families as well as on
pixel AUC ($97.8\%$ against $96.2\%$). Fig.~\ref{fig:qual}
shows the maps themselves.

\begin{table}[t]
\centering
\caption{Spoof-evidence localisation, IoU (\%) against the generator's own mask,
protocol P1, $12$-epoch recipe, full test set, maps binarised at $0.5$. Pixel AUC is
over $40$ samples per family. ResNet-18 and ViT-Tiny have no localisation output.}
\label{tab:loc}
\small
\resizebox{\linewidth}{!}{\input{tables/localisation}}
\end{table}

Predicted maps are binarised at $0.5$, the midpoint of a sigmoid trained with a BCE
plus Dice objective. To check that this convention does not favour either model, we
also swept the cut from $0.10$ to $0.90$ on the validation split. The best cut on validation is $0.55$ for both models, and at that cut the test
region-level IoU moves by less than $0.4$ points for either ($80.1\%$ and $73.5\%$),
so the comparison does not depend on the choice.

\subsection{Is the explanation faithful?}
An explanation is only worth reading if perturbing what it highlights changes the
decision. Table~\ref{tab:faith} reports deletion and insertion AUC
\cite{petsiuk2018rise} of the model's own attack score for the localisation map,
the channel quoted in the generated text, against a random control. The control is
a smoothed random field at the same spatial scale as the map, not per-pixel noise,
which would destroy high-frequency texture everywhere and make the control look
artificially strong. We use $48$ test samples per attack family and give $95\%$
bootstrap intervals for the per-sample difference between map and control.

\begin{table}[t]
\centering
\caption{Faithfulness of the X-UAD localisation map against a smoothed random
control, protocol P1, $12$-epoch model. Insertion gain is map minus control; an
interval below zero means the map does worse than random.}
\label{tab:faith}
\footnotesize
\resizebox{\linewidth}{!}{\input{tables/faithfulness}}
\end{table}

Deletion is clean. Removing what the map highlights lowers the attack score
faster than removing random regions, on the region-level families ($0.892$ against
$0.941$) and on the whole-frame families ($0.904$ against $0.943$), and both
intervals for the gain lie well above zero. The pixels the map points at do carry
the decision.

Insertion is not clean. Restoring the highlighted pixels into a blurred image
recovers the score more slowly than restoring random ones, and the interval is
below zero in both groups. A natural explanation is that the failure comes from the
whole-frame families, where a small highlighted region cannot hold evidence spread
over the whole capture. The data support that only in part: the
shortfall is three times larger on whole-frame families ($-0.021$) than on
region-level ones ($-0.007$), but it does not vanish on the latter. Anyone relying
on these maps should know that they pass the deletion test and fail the insertion
test.

\subsection{White-box attacks on the detector}\label{sec:whitebox}
The adversarial family in our benchmark is a transfer attack crafted against a
separate surrogate. Here we give the attacker the detector itself. We run
$\ell_\infty$ PGD ($10$ steps, step size $\varepsilon/4$, random start) on $168$ test
attacks, $24$ per family, with two goals. The first is evasion: push an attack
towards the genuine class. The second targets X-UAD alone and asks whether its
vacuity can be lowered independently of the decision: we minimise $u$ while a
cross-entropy term holds the predicted class fixed. Table~\ref{tab:whitebox} gives
the results.

\begin{table}[t]
\centering
\caption{White-box PGD on $168$ test attacks, protocol P1, $12$-epoch models, at the
validation threshold. Upper rows: share of attacks still detected under evasion.
Lower rows: X-UAD's mean vacuity under the false-confidence attack, and how often
its predicted class is unchanged.}
\label{tab:whitebox}
\footnotesize
\resizebox{\linewidth}{!}{\input{tables/whitebox}}
\end{table}

Neither detector survives evasion. At $\varepsilon=1/255$ X-UAD still catches
$30.4\%$ of the attacks against $5.4\%$ for CDCN++, but at $2/255$ both are near
$10$--$15\%$, and from $4/255$ X-UAD detects none. CDCN++ levels off near $40\%$ at
the larger budgets, which more likely reflects gradients that are harder to follow
than real robustness; we did not run gradient-free attacks to check. X-UAD's
vacuity does not catch what the classifier misses: none of the evaded attacks
reaches the $u\ge0.55$ referral threshold at any budget.

The second attack shows that the uncertainty channel can be moved on its own. At
$\varepsilon=1/255$ it cuts mean vacuity from $0.064$ to $0.017$ while $99.4\%$ of
predicted classes stay the same. At larger budgets this targeted attack no longer
holds the class fixed (about a third to a half of predictions are kept), so the
clean decoupling is a small-perturbation effect, but that is the regime an attacker
would use.

\subsection{Ablations}\label{sec:abl}
Table~\ref{tab:abl} removes one component at a time at a reduced training budget so
the sweep stays affordable. These rows are comparable with each other and not with
Table~\ref{tab:main}, and we would ask readers not to cross-reference the two.

\begin{table*}[t]
\centering
\caption{Component ablation at a reduced training budget. Deltas are relative to the
full model in the first row. These rows are comparable with each other but not
directly with Table~\ref{tab:main}.}
\label{tab:abl}
\small
\input{tables/ablation}
\end{table*}

The central-difference stem is clearly load-bearing. Removing it costs more than
every other component combined, dropping eight-way accuracy by $21$ points and
region-level IoU by $12$. Localisation supervision is what gives the
evidence map meaning; without it, region-level IoU falls by more than half even
though raw detection accuracy barely moves, which shows
that a saliency map obtained for free from a detector is not the same thing as
a supervised evidence map.

Two rows trade closed-set accuracy for
the uncertainty channel. Removing PUES \emph{improves} reduced-budget ACER,
while removing the capability the term was built to provide.
Table~\ref{tab:oablate} isolates that trade on the fold where it matters,
the face-swap leave-one-attack-out protocol, where PUES is worth $5.4$ novelty AUROC
points and $5.2$ points of family accuracy, at a cost of $2.2$ ACER points. Adaptive
Quantum Attention shows the same pattern. At reduced budget it costs ACER but buys
calibration and localisation, and specifically on the leave-one-attack-out fold it
buys $1.9$ AUC and $1.9$ accuracy at equal novelty AUROC.

\begin{table}[t]
\centering
\caption{Open-set ablation on the face-swap leave-one-attack-out fold, where the
uncertainty channel is exercised.}
\label{tab:oablate}
\small
\resizebox{\linewidth}{!}{\input{tables/oablate}}
\end{table}

\subsection{Efficiency and the Born-channel gate}
X-UAD adds its extra capability without a large parameter cost
(Table~\ref{tab:eff}). The rotation-circuit parameterisation costs
$L\lfloor d/2\rfloor$ angles per circuit against the $d^{2}$ of a dense projection,
and the deterministic reasoning module runs in under a millisecond and adds zero
learned parameters. At $1.29$M parameters, X-UAD sits between ViT-Tiny and
ResNet-18 in size while being the only model in the comparison that outputs a
localisation map, a set of named evidence modes and a written rationale. GMAC counts
exclude attention matmuls, so X-UAD's figure is understated relative to the pure
CNNs.

\begin{table}[t]
\centering
\caption{Inference cost, CPU, single $112\times112$ crop.}
\label{tab:eff}
\small
\resizebox{\linewidth}{!}{\input{tables/efficiency}}
\end{table}

Table~\ref{tab:gates} shows where the Born channel is used. The converged
gate opens furthest on print ($0.434$) and adversarial ($0.394$), and sits lowest on
replay ($0.231$). We had expected replay to be the clearest interference family,
given its moir\'e beat, so this is not the pattern we predicted, and we report the
measurement rather than the prediction.

\begin{table*}[t]
\centering
\caption{Converged adaptive gate $g$, the fraction of the attention score taken from
the Born channel, averaged over heads, layers, and test samples of each family.}
\label{tab:gates}
\small
\input{tables/gates}
\end{table*}

\subsection{Qualitative behaviour}\label{sec:qual}
Fig.~\ref{fig:qual} shows, for one sample per family, the input, the predicted
evidence map, the ground-truth mask and the attention rollout. Table~\ref{tab:modes}
lists the quantum evidence modes with the names read off the trained model, and
Fig.~\ref{fig:report} reproduces a verdict verbatim from \texttt{scripts/demo.py}.
These analyses, and the gate values of Table~\ref{tab:gates}, use the $10$-epoch
original-recipe P1 model.

\begin{table}[t]
\centering
\caption{The quantum evidence basis after training. Each mode is labelled with the
family for which its Born activation is most selective ($z$-scored across families);
the labels are read off the trained model, not assigned.}
\label{tab:modes}
\small
\resizebox{\linewidth}{!}{\input{tables/modes}}
\end{table}

\begin{figure}[t]
\centering
\includegraphics[width=\linewidth]{figures/qualitative.png}
\caption{One sample per family. Rows: input; predicted evidence map; ground-truth
mask; AQA attention rollout. The localiser recovers the eye/mouth support of the
reenactment, the cut-outs of the 3-D mask, and spreads over the whole frame for
print, synthetic and adversarial. It is not perfect: on this face-swap it predicts
an eye/mouth pattern resembling reenactment rather than the ground-truth blending
ellipse, and on replay it misses the localised glare spot.}
\label{fig:qual}
\end{figure}

\begin{figure}[t]
\centering
\fbox{\parbox{0.95\linewidth}{\ttfamily\footnotesize
\input{tables/sample_report}
}}
\caption{A verdict as printed by the deployed system. Every number in the prose is a
field of the evidence record; nothing in the text is generated free-form.}
\label{fig:report}
\end{figure}

% =====================================================================
\section{The Research Gap, Revisited}\label{sec:gap}
% =====================================================================
We restate the gap from Section~\ref{sec:related} as five concrete deficits, and pair
each one with the mechanism that addresses it and the measurement that shows the
effect.

\textbf{Gap 1: detection is unified only in-distribution.} The unified corpora
\cite{yu2024unified,ma2025uniattackplus} show near-saturated seen-attack numbers next
to $22$ to $44\%$ ACER once attack families are withheld. Nothing in those
architectures represents ``this belongs to no family I know'', because a softmax has
to put all of its mass somewhere. \emph{Mechanism:} the Dirichlet evidential head
makes total evidence, and therefore vacuity, an explicit output, and PUES trains that
output against an outlier proxy built from the attack taxonomy itself rather than an
unrelated image pool. \emph{Evidence:} the unseen-family detection rate and novelty
AUROC in Table~\ref{tab:loao}, and the w/o-PUES row of Table~\ref{tab:oablate}.

\textbf{Gap 2: attention is linear, but much of the evidence is interference.}
Moir\'e beats, halftone screens, and upsampling checkerboards are periodic structures
whose sign carries no information. A dot-product score is sign-sensitive and has to
spend capacity learning that invariance from scratch. \emph{Mechanism:} the Born
channel scores token pairs by a squared overlap, which is sign-invariant by
construction, and the input-conditioned gate recruits it only where it helps.
\emph{Evidence:} the w/o-AQA rows of Tables~\ref{tab:abl} and~\ref{tab:oablate}, and
the per-family gate values in Table~\ref{tab:gates}. We should be careful about what
the gate table does and does not show. It confirms the channel is used,
with a mean around $0.31$ and a clear per-family structure, and the ablation confirms
it contributes. It does \emph{not} confirm our original hypothesis that the gate
would open furthest on the most periodic families. Print and adversarial open widest,
while replay, the family with the most obvious moir\'e beat, opens least. The
mechanism works; our story about why it works is not settled by this measurement.

\textbf{Gap 3: pixel supervision targets a liveness proxy, not the evidence itself.}
Depth maps \cite{liu2018auxiliary,yu2021revisiting} are undefined for a digitally
manipulated but physically genuine capture, so a unified detector built on them
inherits a supervision signal that is meaningless for half of its own label space.
\emph{Mechanism:} we supervise the actual injected support directly, which our
benchmark can provide because it knows exactly where it wrote the evidence.
\emph{Evidence:} Table~\ref{tab:loc}'s region-level IoU on families for which a depth
target would carry zero information, and Fig.~\ref{fig:qual}.

\textbf{Gap 4: MLLM explanations are post-hoc narration.} FFAA \cite{huang2024ffaa}
and M2F2-Det \cite{guo2025m2f2det} give the language model the raw image, not the
detector's internal state, so the narrative is unverifiable and can contradict the
decision it is meant to explain. \emph{Mechanism:} we invert the dependency. The
forward pass emits a typed evidence record, and language is realised from that record
alone, with the optional hosted back-end constrained and then re-verified against the
same record. \emph{Evidence:} the grounded realiser's hallucination rate is zero by
construction, since every sentence is gated on a field of the record
(Fig.~\ref{fig:report}). Whether the localisation map it quotes is itself faithful
is a separate question, and Table~\ref{tab:faith} gives a mixed answer: removing the highlighted pixels changes the decision more than removing random ones,
but restoring them does not recover it faster than random restoration.

\textbf{Gap 5: a score is not a decision.} An operator has to act on something, and a
scalar in $[0,1]$ with no calibrated meaning and no severity model does not tell them
what to do. \emph{Mechanism:} the risk calibrator converts posterior, vacuity and
evidence area into a bounded risk score with an explicit unknown-attack premium, so
that a confidently-unknown sample is escalated to human review instead of silently
passing through. \emph{Evidence:} the verdict of Fig.~\ref{fig:report} shows the
mechanism working on a single sample. We do \emph{not} yet have evidence that the
$u\ge0.55$ referral rule helps on unseen attacks. In every leave-one-attack-out
fold, genuine faces receive at least as much mean vacuity as the held-out family
($0.54$, $0.37$, $0.42$, $0.58$ against $0.54$, $0.31$, $0.43$, $0.39$), so the rule
as set is likely to refer at least as many genuine faces as unseen attacks. Setting
the referral threshold, and measuring its referral rates per group, is left to
future work.

% =====================================================================
\section{Limitations}\label{sec:limits}
% =====================================================================

All numbers come from a single seed on a simulated benchmark, and our adapters for
seven public corpora have not yet been run on captured data. X-UAD received more
tuning than the baselines, and the P2, P3 and ablation tables come from an earlier
code version whose numbers may not reproduce exactly. The models are small,
image-only and undefended against white-box attacks.

% =====================================================================
\section{Open Directions}\label{sec:open}
% =====================================================================
A few of the questions below go beyond what this paper set out to do, and we name
them explicitly rather than folding them into the limitations list, because we are
not aware of existing work in this space that has tackled them yet.

\textbf{Defending the uncertainty channel.} Open-set and evidential work, ours
included, usually tests the robustness of the class decision and not of the
uncertainty signal. Section~\ref{sec:whitebox} shows that the gap is real: with
white-box access, a perturbation of $1/255$ cuts X-UAD's mean vacuity by about
three quarters while leaving $99.4\%$ of its predicted classes unchanged. Our vacuity head reads a
stop-gradient copy of the shared representation, so the two objectives are
decoupled by design, which is what makes the uncertainty channel free with respect
to accuracy and probably also what makes it so easy to move on its own. We do not
know of a defence designed for this case. Adversarial training on the vacuity
objective, or tying $u$ to an input-space consistency check, are the obvious first
things to try.

\textbf{PUES beyond pairwise blends.} We only ever mix two attack families at a time.
Whether three-way or higher-order mixtures, or mixtures weighted by how forensically
\emph{similar} two families already are rather than a flat $\mathrm{Beta}(2,2)$,
produce a better outlier proxy is untested. There is a plausible argument either way.
Closer families might produce a harder and more informative training signal for the
vacuity boundary, or they might produce blends that are accidentally too close to a
real attack and corrupt the class-direction gradient.

\textbf{A learned, rather than fixed, risk policy.} Section~\ref{sec:method}
deliberately keeps the risk calibrator's coefficients unlearned, for auditability
reasons we still believe in. But nothing stops a clearly-labelled meta-model from
learning site-specific coefficients from a small amount of labelled operator
feedback, such as analyst overrides of the band, while leaving the underlying
detector and its explanation completely untouched. That would let a deployment site's
risk tolerance adapt over time without ever touching the part of the system
responsible for the actual forensic decision.

\textbf{Cross-domain shift and attack-family shift as a joint problem.} Every
domain-generalisation technique for face anti-spoofing we are aware of
\cite{wang2022ssan,shao2019mada,jia2020ssdg} targets capture-domain shift alone,
using only \emph{seen} attack families. Every open-set technique, including PUES,
targets attack-family shift on a \emph{single} capture domain. What happens at the
intersection, where an unseen attack family arrives through an unseen capture domain
simultaneously, has not been studied. Our own Table~\ref{tab:cross} suggests the joint
case may not be the product of the two individual failure rates. Closing this
gap would need a benchmark explicitly designed to vary both axes independently, which
our generator is well placed to provide.

\textbf{Explanation faithfulness under distribution shift.} Table~\ref{tab:faith}
measures faithfulness only on in-distribution, seen-attack samples. Whether an
intrinsic explanation stays faithful when the underlying detector is operating in the
P2 or P3 regime, which is to say whether the localisation map still points at
something real when the class decision itself is unreliable, is untested here. A
model that is uncertain but still produces a confident-looking localisation map would
be a dangerous failure mode for the auditability reasons this paper
cares about.

% =====================================================================
\section{Conclusion}
% =====================================================================
X-UAD detects physical and digital face attacks with a single model, reports how much
it does not know, shows where the evidence sits, and explains itself in language that
can be audited sentence by sentence. Three ideas make this combination possible: a
Born-rule attention channel that matches the interference character of the underlying
evidence at a parameter cost of $\mathcal{O}(dL)$; an orthonormal quantum evidence
basis that serves simultaneously as a classification feature and as the explanation
itself; and pseudo-unseen evidence suppression, which manufactures an out-of-taxonomy
proxy directly from the attack taxonomy and turns Dirichlet vacuity into a usable
open-set score.

On seen attacks X-UAD matches the strongest baseline on the attack/genuine decision
and localises the evidence better than any model we compared. Its other gains are on
unseen attack families. It also has weaknesses: lower eight-way family accuracy than
both convolutional baselines, no calibration advantage once the baselines are
temperature-scaled, a mildly inverted uncertainty ranking on the adversarial fold,
an insertion-faithfulness result below its control, an uncertainty signal that a
white-box attacker can switch off, and a cross-domain protocol on which no method
here performs acceptably. All numbers come from a single seed on a simulated
benchmark. Repeating them over seeds and validating the mechanisms on captured
corpora are the next steps.

\section*{Acknowledgment}
The source faces used to build the benchmark are drawn from the Labeled Faces in the
Wild database \cite{huang2007lfw}, used under its terms for non-commercial research.

\bibliographystyle{IEEEtran}
\bibliography{refs}

\end{document}
